\section{Conclusion}\label{sec:conclusion}

% Outline:
% - Recap: RADP folds placement and recovery into one DP, with two
%   runtime mechanisms (async mirror cache, async chain) that make
%   that joint optimization viable on a real edge fleet.
% - The strongest empirical claim is the 28-point matrix: latency
%   placement strictly dominates throughput placement across every
%   variant we measured. The underlying R+psi joint optimization is
%   what produces the dominance — not any specific topology or
%   runtime architecture.
% - Async chain forwarding adds 17--47% throughput at C=16 over
%   sync without ever reversing the L>T ordering, confirming the
%   two mechanisms are *orthogonal*: chain mode is a runtime knob,
%   not a placement input.
% - Future work: deeper sharded models (Llama-2-7B INT4), concurrent
%   fault scenarios (R as a backup-list), cost-function refinement
%   to close the predicted-vs-measured gap.

\todo{Write 3 paragraphs from outline above.}
